% ============================================================
% f2012_guangdong_science_q13_flowchart.tex
% 2012 广东卷（理科）第 13 题 程序框图（TikZ 矢量重绘）
% 开始 → 输入 n → i=2, k=1, s=1 → 判断 i<n
%   是：s=(1/k)(s×i) → i=i+2 → k=k+1 → 回到判断
%   否：输出 s → 结束
% ============================================================
\begin{tikzpicture}[
  x=1cm, y=1cm, line width=0.6pt, font=\footnotesize, >=Stealth,
  every node/.style={inner sep=3pt},
  proc/.style={rectangle, draw, minimum width=2.9cm, minimum height=0.68cm},
  io/.style={trapezium, draw, trapezium left angle=70, trapezium right angle=110,
             minimum width=3.0cm, minimum height=0.72cm, trapezium stretches=true},
  dec/.style={diamond, draw, aspect=2.3, minimum width=2.5cm, minimum height=0.95cm},
  term/.style={rectangle, rounded corners=8pt, draw, minimum width=1.9cm, minimum height=0.60cm}
]

  \node[term] (start) at (0,0) {开始};
  \node[io]   (input) at (0,-1.25) {输入 $n$};
  \node[proc] (init)  at (0,-2.50) {$i=2$，$k=1$，$s=1$};
  \node[dec]  (judge) at (0,-3.95) {$i<n$};

  \node[proc, minimum height=1.10cm] (body) at (0,-5.55) {$s=\dfrac1k(s\times i)$};
  \node[proc] (inc)   at (0,-6.95) {$i=i+2$};
  \node[proc] (kk)    at (0,-8.20) {$k=k+1$};

  \node[io]   (out)   at (3.90,-5.55) {输出 $s$};
  \node[term] (end)   at (3.90,-6.95) {结束};

  \draw[->] (start) -- (input);
  \draw[->] (input) -- (init);
  \draw[->] (init) -- (judge);
  \draw[->] (judge) -- node[right] {是} (body);
  \draw[->] (body) -- (inc);
  \draw[->] (inc) -- (kk);
  \draw[->] (kk.west) -- (-2.55,-8.20) -- (-2.55,-3.95) -- (judge.west);
  \draw[->] (judge.east) -- (3.90,-3.95) -- (out.north);
  \node at (2.55,-3.72) {否};
  \draw[->] (out) -- (end);

\end{tikzpicture}
